% A284 -- Messung, Projektion und der Ursprung des Probabilismus
% Johann Pascher -- ORCID 0009-0000-6518-4064
% Kapitelquelle zur A-Serie; Wrapper: wr_standalone_A4/A284_Messung_Projektion_De.tex
% Status: August 2026

\providecommand{\qedsq}{\raisebox{0.03em}{\framebox[0.62em]{\rule{0pt}{0.42em}}}}

% ===================== KURZFASSUNG =====================
\begingroup\small\noindent
\textbf{Kurzfassung.}
Dieses Dokument schließt eine Lücke, die in A070, A090 und A160 jeweils
nur angedeutet ist: die geschlossene Argumentation dafür, \emph{warum}
Messung in der FFGFT das Gesamtsystem beeinflusst, und \emph{warum}
die Auswertung trotz deterministischer Grundgleichungen probabilistisch
ist. Die drei Ebenen werden getrennt geführt und am Ende zur
Instantanitätsfrage der Bell-Korrelationen zusammengeführt.

Kernthesen:
\textbf{(1)} Messung ist eine geometrische Typ-II-Projektion (A090) --- sie
wirkt auf den globalen Zustand auf $T^4$, nicht auf einen lokalen Teilbereich.
\textbf{(2)} Der Probabilismus entsteht nicht aus fundamentalem Zufall, sondern
aus dem Informationsverlust der Projektion: die Born-Regel ist die Statistik des
Typ-II-Verlusts.
\textbf{(3)} Was als instantane Fernwirkung zwischen zwei Messpunkten A und B
erscheint, ist die statistische Signatur einer globalen Mode, abgelesen an zwei
Projektionsstellen desselben Systems --- kein Prozess läuft von A nach B.
\par\endgroup

\bigskip

\begingroup\small\noindent
\textbf{Schichtmarkierungen.}
\textbf{[SETZUNG]} deklarierter Ausgangspunkt ---
\textbf{[B]} algebraisch/mathematisch abgeleitet ---
\textbf{[K]} numerisch aus Korpus-Konstanten berechnet ---
\textbf{[S]} Plausibilitätsskizze ---
\textbf{[H]} offene Forschungsfrage.
\par\endgroup

\bigskip

% ===================== 1 =====================
\section{Worum es geht}

\textbf{[SETZUNG]} In der Standardquantenmechanik ist der Messprozess
das konzeptuell unbefriedigendste Element: eine einheitliche, deterministische
Evolution der Wellenfunktion bricht beim Messen zusammen, ein probabilistisches
Ergebnis erscheint, und zwei weit entfernte Teilchen scheinen instantan
voneinander zu wissen. Drei Jahrzehnte nach der Bell-Verletzung bleibt die Frage
offen, was dabei physikalisch passiert.

\textbf{[B]} Die FFGFT hat die Elemente einer Antwort bereits in anderen
Dokumenten verteilt: A090 (Projektionskette, Typ-I/II/III), A070
(Messung als Polübergang, Born-Regel aus Phasenunkenntnis), A160
(Bell-Korrelationen als topologische Verbindung), A030 (Windungsklassen
ohne Ort). Was fehlt, ist die geschlossene Argumentationskette.
Dieses Dokument zieht sie.


% ===================== NEUER ABSCHNITT =====================
\section{Lichtstrahlen als Beobachterabstraktion --- das EM-Feld als Träger}
\label{sec:lichtstrahlen}

\subsection{Was ein Lichtstrahl wirklich ist}

\textbf{[S]} Ein Lichtstrahl ist keine physikalische Größe. Er ist eine
\emph{geometrische Abstraktion}, die aus dem EM-Feld unter bestimmten
Bedingungen destilliert wird --- nämlich genau dann, wenn die Wellenlänge
$\lambda$ klein ist gegen alle relevanten geometrischen Abmessungen
(Blenden, Objekte, Beobachterabstand). Diese Bedingung ist die
\emph{Strahlenoptik als Grenzfall der Wellenoptik} --- mathematisch:
$\lambda \to 0$ in der Wellengleichung ergibt die Eikonal-Gleichung,
aus der Strahlen als Normalentrajektorien der Phasenflächen folgen.

\textbf{[B]} Das EM-Feld ist der Träger. Es existiert und entwickelt
sich nach den Maxwell-Gleichungen, vollständig beobachterunabhängig.
Der Lichtstrahl erscheint erst, wenn ein Beobachter oder eine Detektoranordnung
eine Ortsmessung macht: Der Detektor reagiert auf das Feld an seinem Ort
und \emph{konstruiert} rückwärts einen Strahl als die geometrische Linie,
die Quelle und Detektor verbindet. Der Strahl ist das Ergebnis dieser
Konstruktion --- nicht sein Vorläufer.

\textbf{[B]} Dieser Sachverhalt ist nicht neu --- er liegt in der
klassischen Wellenoptik vollständig vor. Was in der FFGFT hinzukommt:
Die Feldentwicklung auf $T^4$ hat keine Orte im Standardsinn. Ein Photon
hat keine Trajektorie; es hat eine Modenstruktur auf dem Torus. Was wir
als ``Lichtweg'' beobachten, ist die statistische Häufung von
Detektionsereignissen in einer Raumregion --- eine Projektionststatistik,
keine physikalische Bahn (A090, Typ-II-Projektion).

\subsection{Warum es so schwer ist, sich von der Lichtwegvorstellung zu lösen}

\textbf{[S]} Die Lichtwegvorstellung sitzt aus vier Gründen tief:

\textbf{(1) Evolutionär.} Das menschliche Gehirn ist für Kausalität
durch räumliche Übertragung optimiert: Stein fliegt, trifft, wirkt.
``Ursache vor Wirkung, Träger verbindet, Ferne trennt'' --- diese
Heuristiken waren $10^5$ Jahre lang überlebensrelevant. Sie sitzen
tiefer als jede Theorie.

\textbf{(2) Sprachlich.} Jede natürliche Sprache hat Subjekte, die
Prädikate ausführen: ``Das Licht geht von A nach B.'' Die Idee, dass
ein Feld existiert und ein Beobachter lokal reagiert --- ohne dass
``etwas'' von A nach B läuft --- ist grammatisch nicht natürlich
formulierbar. Selbst physikalisch korrekte Sätze klingen wie
Verneinungen, nicht wie positive Aussagen.

\textbf{(3) Historisch.} Newton dachte noch in Korpuskeln --- Lichtteilchen
auf Bahnen. Maxwell löste das: Felder statt Bahnen. Aber dann kamen
Photonen, und die Alltagsphysik verfiel zurück in die Bahnvorstellung:
``Das Photon geht durch den Doppelspalt.'' Die Wellenoptik wurde zur
Rechenmethode, die Teilchenbahn zur Anschauung. Dieser Rückschritt ist
tief in die Physikausbildung eingebaut.

\textbf{(4) Sichtbarkeit durch Umwege.} Lichtstrahlen werden nur über
Umwege sichtbar: Staub, Rauch, Tröpfchen --- Streuzentren, die das
Feld lokal absorbieren und reemittieren. Was wir im Scheinwerfer-Strahl
sehen, sind Millionen von Streuereignissen, jedes an einem anderen Ort.
Der ``Strahl'' ist die statistische Linie dieser Ereignisse. Ohne
Streuzentren ist das EM-Feld unsichtbar --- physikalisch real, aber
kein Beobachtungsobjekt.

\subsection{Warum wirken dann Linsen?}

\textbf{[B]} Die Linsenfrage ist der präziseste Test der These.
Wenn Lichtstrahlen Abstraktionen sind und das EM-Feld der Träger
ist --- warum fokussiert dann eine Linse?

Die Antwort liegt in zwei Schichten:

\textbf{Schicht 1 --- Wellenoptik (vollständige Beschreibung):}
Das EM-Feld trifft auf das Glasmedium der Linse. Im Glas ist die
Phasengeschwindigkeit $v = c/n$ kleiner als im Vakuum. Die
Wellenfront --- eine Fläche gleicher Phase --- verlangsamt sich im
dickeren Teil der Linse stärker als im dünneren Teil. Nach dem
Durchgang ist die Wellenfront nicht mehr eben, sondern gekrümmt:
Sie konvergiert. Das Feld dahinter entwickelt sich auf eine
Fokalpunkt-Struktur zu --- eine Region maximaler Feldintensität.
Kein Strahl ist nötig: Die Maxwell-Gleichungen im inhomogenen
Medium beschreiben alles vollständig.

\textbf{Schicht 2 --- Strahlenoptik (Grenzfall-Beschreibung):}
Im Grenzfall $\lambda \ll f$ (Wellenlänge $\ll$ Brennweite) gilt
Snellius: $n_1 \sin\theta_1 = n_2 \sin\theta_2$. Die Strahlen sind
jetzt Normalentrajektorien der Wellenfronten --- sie beschreiben
\emph{wohin das Feld sich konzentriert}, nicht woher ein Teilchen
kommt. Die Linse lenkt keine Körper; sie formt Phasenflächen um.

\textbf{[B]} Das Entscheidende: Die Linse wirkt auf das Feld, nicht
auf Strahlen. Der Strahl ist eine nützliche geometrische Beschreibung
des Ergebnisses --- aber das Feld ist die Ursache. Ein Beobachter
hinter der Linse sieht das Bild, weil das Feld an seinem Detektorort
eine hohe Intensität hat. Er \emph{konstruiert} daraus einen Strahlengang.

\textbf{[S]} In der FFGFT-Lesart: Das Photon ist eine Feldmode auf
$T^4$. Es hat keine Bahn. Die Linse modifiziert die lokale Metrik
des EM-Feldes --- sie verändert die Phasenstruktur der Mode. Hinter
der Linse ist die Wahrscheinlichkeit, das Photon am Fokalpunkt zu
detektieren, maximal. Der Fokalpunkt ist keine Station auf einer Bahn,
sondern ein Maximum der Projektionswahrscheinlichkeit. Dass Linsen
``funktionieren'', ist eine Aussage über die Verteilung von
Typ-II-Projektionsereignissen im Raum --- nicht über Bahnen.

\subsection{Die Verbindung zur Bell-Frage}

\textbf{[B]} Der Zusammenhang mit der Instantanitätsfrage ist direkt:
Wer fragt ``Wie kann die Messung an A instantan B beeinflussen?''
denkt in Lichtwegen. Der Lichtweg ist die implizite Annahme, dass
jede Wirkung einen räumlichen Träger hat, der sich von A nach B bewegt.

Aber wenn das EM-Feld der Träger ist --- und das Feld global auf $T^4$
existiert --- dann gibt es kein ``von A nach B''. Das Feld ist schon überall.
Die Messung an A ist eine lokale Wechselwirkung mit dem globalen Feld.
Die Korrelation an B ist die Antwort desselben globalen Feldes an einem
anderen Projektionspunkt. Kein Träger muss reisen.

\textbf{[S]} Das Linsenbeispiel macht das greifbar: Wenn Licht durch
eine Linse geht und am Fokalpunkt landet, fragt niemand: ``Welcher
Strahl hat den Fokalpunkt verursacht?'' Das ist die falsche Frage.
Das Feld hat die Phasenstruktur, die zu diesem Fokalpunkt führt.
Genauso: Wenn ein verschränktes Photonenpaar gemessen wird, fragt
die FFGFT nicht: ``Welches Signal lief von A nach B?'' Das Feld
hat die Modenstruktur, die zu diesen Korrelationen führt.


% ===================== 2 =====================
\section{Was Messung in der FFGFT ist}

\subsection{Messung ist eine Typ-II-Projektion}

\textbf{[B]} In der Projektionskette (A090) gibt es drei qualitativ
verschiedene Operationen: Typ~I (Dualitäts-Dekompaktifizierung, fast
reversibel), Typ~II (geometrische Projektion, verlustbehaftet, nicht
reversibel), Typ~III (repräsentationelle Übersetzung, bijektiv).

\textbf{[B]} Eine Messung ist eine Typ-II-Operation. Sie kollabiert eine
Dimension des Zustandsraums --- sie projiziert den globalen Zustand auf
$T^4$ auf einen niederdimensionalen Unterraum. Der Kern dieser Projektion
$\ker(\pi) = \mathbb{R}$ enthält Information, die nach der Projektion nicht
mehr im Bild enthalten ist (A090). Diese verlorene Information ist nicht weg ---
sie ist im Messapparat deponiert, der damit Teil des Gesamtsystems wird.

\textbf{[B]} Drei Konsequenzen folgen unmittelbar:

\begin{enumerate}[leftmargin=1.6em,itemsep=0.3em]
  \item \textbf{Messung wirkt global.} Der $T^4$ ist randlos
    ($\partial T^4 = \emptyset$, A030) und trägt globale Eigenmoden ohne
    ausgezeichneten Ort. Eine Projektion an einem Projektionspunkt ändert
    den globalen Modeninhalt --- es gibt keine lokale Messung, die den Rest
    des Systems unberührt ließe.
  \item \textbf{Der Messapparat ist Teil des Systems.} Was vorher
    Systemzustand + Messapparat-Zustand war, ist nachher ein gemeinsamer
    Zustand. Die Grenze ist keine physikalische Grenze, sondern eine
    Modellgrenze (Schnittstellendefinition).
  \item \textbf{Wiederholung verändert das System.} Jede weitere Messung
    trifft auf einen anderen globalen Zustand als die erste. Der Torus
    \enquote{weiß} von der ersten Messung --- nicht durch ein Signal,
    sondern weil sein globaler Modeninhalt sich verändert hat.
\end{enumerate}

\subsection{Die geometrische Lesart des Kollapses}

\textbf{[S]} In der Standard-QM kollabiert die Wellenfunktion beim Messen
auf einen Eigenzustand. In der FFGFT-Lesart (A070) ist das eine
deterministische Bewegung im Modenraum: der $(z,r,\theta)$-Punkt wird auf
das nächstgelegene stabile Extrem getrieben ($z=+1$ oder $z=-1$). Das ist
keine instantane Aktion auf Distanz, sondern eine lokale geometrische
Bewegung am Projektionspunkt. Der \enquote{Kollaps} ist der
Typ-II-Projektionsschritt selbst.

\textbf{[B]} Was dabei verloren geht, ist die Phase $\theta$ ---
die azimutale Information der Mode. Sie ist nach der Messung im Bild
nicht mehr zugänglich, aber nicht vernichtet: sie ist im Messapparat
kodiert (Typ-II-Verlust).

% ===================== 3 =====================
\section{Woher der Probabilismus kommt}

\subsection{Drei mögliche Erklärungen}

Es gibt drei strukturell verschiedene Antworten auf die Frage, warum
quantenmechanische Messungen probabilistische Ergebnisse liefern:

\textbf{Standardlesart (ontologischer Zufall):} $|\psi|^2$ ist eine
ontologische Wahrscheinlichkeit. Der Ausgang ist fundamental zufällig,
nicht durch verborgene Variablen bestimmt (Bell, Kochen-Specker).

\textbf{Verborgene-Variablen-Lesart (lokal):} Der Ausgang ist durch
verborgene lokale Parameter bestimmt. Diese Lesart ist durch Bell
experimentell ausgeschlossen.

\textbf{FFGFT-Lesart (Projektionsverlust):} Der Ausgang ist durch
die deterministischen Feldgleichungen auf $T^4$ bestimmt, aber
die Phase $\theta$ ist für den Beobachter nicht zugänglich.
Die Stochastizität entsteht beim Messapparat (der Typ-II-Projektion),
nicht im Naturprozess.

\subsection{Warum die FFGFT-Lesart Bells Theorem nicht verletzt}

\textbf{[B]} Bell schließt \emph{lokale} verborgene Variablen auf dem
projizierten $\mathbb{R}^3$ aus. Die Phase $\theta$ auf dem $T^4$
ist keine lokale verborgene Variable im Sinne des Theorems --- sie ist
ein globales Merkmal des Torusmodes, kein Parameter der an einem
Raumort sitzt. Bells Beweis setzt eine Faktorisierung der
verborgenen-Variablen-Dichte über die räumliche Trennung voraus.
Diese Faktorisierung gilt auf $T^4$ nicht, weil Orte auf $T^4$
durch ihre Windungsklassen definiert sind, nicht durch Raumkoordinaten
(A030).

\textbf{[B]} Die quantitative Form: Die Born-Regel
\[
  P_\text{FFGFT}(z=+1) = \frac{1+z}{2} = |\alpha|^2 = P_\text{QM}(0)
\]
folgt aus der Gleichverteilung über $\theta \in [0,2\pi)$ bei großer
Statistik (A070). Der Probabilismus ist das Ensemble-Mittel über alle
Phasenzustände --- nicht fundamentaler Zufall, sondern Mittelung über
das, was der eingebettete Beobachter nicht sehen kann.

\subsection{Born-Regel als Projektionststatistik}

\textbf{[S]} Die strukturelle These: Die Born-Regel $P = |\psi|^2$ ist
die Statistik des Typ-II-Informationsverlusts. Bei einer Projektion
$\pi: \mathbb{C}^n \to \mathbb{C}^{n-1}$ verteilt sich die Information
der verlorenen Dimension auf das Ensemble der Messwiederholungen.
Die quadratische Form entsteht, weil die Norm im Hilbertraum euklidisch
ist und die Projektion isometrisch auf den Unterraum wirkt.

\textbf{[H]} Die quantitative Verbindung zwischen dem Typ-II-Verlust
und der exakten Form der Born-Regel ($|\psi|^2$ statt $|\psi|$ oder
$|\psi|^3$) ist nicht ausgearbeitet. Das ist die eigentliche offene
Stelle --- nicht ob, sondern wie genau die Projektionsstatistik die
Quadratform erzwingt.

% ===================== 4 =====================
\section{Drei Ebenen der Nichtlokalität}

\textbf{[B]} Die Instantanitätsfrage bei Bell-Tests hat drei Ebenen,
die strikt getrennt werden müssen. Alle drei sind im FFGFT-Rahmen
konsistent auflösbar --- auf verschiedene Weise.

\subsection{Ebene 1 --- Torusbild: topologische Nachbarschaft}

\textbf{[SETZUNG]} (A160) Bell-Korrelationen sind auf dem $T^4$ eine
topologische Erscheinung. Was im projizierten $\mathbb{R}^3$ als
Fernwirkung erscheint, sind im kompakten Bild durch geschlossene Wege
verbundene Punkte --- topologische Nachbarn.

\textbf{[B]} Für irdische Abstände $d \ll L^*$ gilt: Beide Messpunkte
A und B liegen auf demselben Torusblatt. Der $T^4$ mit der
Abschlussskala $L^* = 4\bar\lambda_e/\xi_0^{10} \approx 1{,}3 \times
10^{26}$\,m ist auf irdischen Skalen ($d \sim 10^6$\,m) praktisch flach.
\emph{Aber:} Die topologische Verbindung läuft nicht über die euklidische
Distanz, sondern über die Windungsklasse. Ein Verschränkungspaar teilt
dieselbe Windungsklasse --- das ist der Sinn der topologischen
Verbindung, unabhängig vom projizierten Abstand.

\subsection{Ebene 2 --- Feldbild: globale Eigenmode}

\textbf{[B]} (A030) Windungsklassen und Eigenmoden haben keinen Ort.
Eine Eigenmode ist global --- sie ist eine Eigenschaft des gesamten
Torus, nicht eines Bereichs. Ein verschränktes Paar ist kein Paar
getrennter Objekte mit einer Verbindung, sondern \emph{ein} Objekt
mit zwei Projektionsstellen (A160).

\textbf{[B]} Daraus folgt unmittelbar: Es propagiert nichts von A
nach B. Die Korrelation ist Identität, keine Übertragung. Signalfreiheit
ist keine Zusatzbedingung, die man fordern muss --- sie folgt aus der
Konstruktion.

\subsection{Ebene 3 --- Messprozess: globale Projektion}

\textbf{[S]} Das ist die tiefste Ebene, die in den bisherigen
Dokumenten nicht explizit formuliert war.

Die Messung an A ist keine lokale Störung, die danach zu B propagiert.
Sie ist eine Typ-II-Projektion des \emph{globalen} Zustands auf $T^4$.
Diese Projektion verändert den globalen Modeninhalt --- den gemeinsamen
Zustand des Gesamtsystems, zu dem beide Projektionsstellen A und B
gehören. Was an B als \enquote{instantane Auswirkung} erscheint, ist
keine Wirkung die von A kommt, sondern die veränderte Randbedingung
des globalen Zustands, aus dem B abliest.

\textbf{[S]} Anders formuliert: Die Messung an A liest die globale
Mode ab und modifiziert sie dabei. Die nächste Messung an B trifft
auf eine bereits modifizierte globale Mode --- nicht weil ein Signal
von A nach B lief, sondern weil beide Teil desselben globalen Systems
sind.

\subsection{Warum die Auswertung immer probabilistisch bleibt}

\textbf{[B]} Alle drei Ebenen konvergieren auf denselben Punkt:
Die Auswertung ist probabilistisch, weil:

\begin{enumerate}[leftmargin=1.6em,itemsep=0.3em]
  \item Jede Einzelmessung eine Typ-II-Projektion ist (Informationsverlust
    der Phase $\theta$).
  \item Die globale Mode nach jeder Messung modifiziert ist ---
    kein zweites Mal ist der Ausgangszustand exakt derselbe.
  \item Das Ensemble über viele Messungen die Phase-Mittelung liefert,
    aus der die Born-Regel folgt.
\end{enumerate}

\textbf{[B]} Man kann keine Bell-Korrelation an einer Einzelmessung
ablesen. Man braucht immer ein Ensemble. Das Ensemble ist die Statistik
der Projektion --- und die Projektion ist global. Deshalb ist die
Auswertung prinzipiell probabilistisch, nicht aus Bequemlichkeit,
sondern weil der Einzelmessprozess die Information systematisch verbirgt.

% ===================== 5 =====================
\section{Feldlaufzeit, nicht Weglaufzeit --- die messbare Verzögerung}

\subsection{Die entscheidende Unterscheidung}

\textbf{[B]} Es gibt zwei grundverschiedene Arten, wie eine Verzögerung
$\Delta t = d/c$ entstehen kann:

\textbf{Weglaufzeit:} Ein Objekt oder ein Signalträger bewegt sich von
Punkt A nach Punkt B. Es gibt ein \emph{Etwas}, das den Weg zurücklegt.
Der Detektor an B registriert das Etwas wenn es ankommt. Diese
Vorstellung ist die intuitive und liegt dem Konzept des Lichtstrahls
zugrunde.

\textbf{Feldlaufzeit:} Das Feld existiert bereits überall. An Punkt A
passiert etwas --- eine Messung, eine Störung, eine Wechselwirkung ---
und diese \emph{Änderung des Feldzustands} breitet sich als Wellenfront
aus. Der Detektor an B registriert nicht ein ankommendes Objekt,
sondern eine veränderte Feldkonfiguration an seinem Ort.

\textbf{[B]} Der mathematische Ausdruck ist für beide Fälle identisch
--- $\delta(|\mathbf{x}-\mathbf{x}'| - c(t-t'))$ --- aber die
physikalische Deutung ist grundverschieden. In der Green-Funktion
(Dok.~131):
\[
  G(\mathbf{x},\mathbf{x}',t-t')
  = \theta(t-t') \cdot \frac{\delta(|\mathbf{x}-\mathbf{x}'| - c(t-t'))}{4\pi|\mathbf{x}-\mathbf{x}'|}
\]
steht $\mathbf{x}'$ für den \emph{Quellpunkt der Störung} und
$\mathbf{x}$ für den \emph{Beobachtungspunkt}. Die $\delta$-Funktion
sagt: Die Änderung erreicht den Beobachtungspunkt genau dann, wenn
$|\mathbf{x}-\mathbf{x}'| = c(t-t')$. Was propagiert, ist die
\emph{Wellenfront der Feldänderung} --- kein Träger.

\subsection{Was bei der Bell-Messung propagiert}

\textbf{[S]} Bei der Messung an Detektor A passiert Folgendes im
Feldterms:

\begin{enumerate}[leftmargin=1.6em,itemsep=0.3em]
  \item Das globale EM-Feld (der verschränkte Zustand) wechselwirkt
    mit dem Messapparat an Ort A. Das ist eine lokale Wechselwirkung
    --- kein globaler Vorgang.
  \item Durch diese Wechselwirkung ändert sich der Feldzustand an A:
    Die Phase, die Amplitude, die Modenbesetzung werden modifiziert.
    Das ist eine \emph{lokale Feldänderung}.
  \item Diese lokale Feldänderung breitet sich als retardierte Welle
    aus --- mit der Green-Funktion $G(\mathbf{x}, \mathbf{x}_A, t-t_A)$.
    Sie erreicht jeden Punkt $\mathbf{x}$ zur Zeit $t_A + |\mathbf{x} -
    \mathbf{x}_A|/c$.
  \item Detektor B, der sich bei $\mathbf{x}_B$ befindet, registriert
    die geänderte Feldkonfiguration zur Zeit $t_A + d_{AB}/c$ --- nicht
    weil ein Signal von A nach B gelaufen ist, sondern weil die
    Feldänderung ihn erreicht hat.
\end{enumerate}

\textbf{[B]} Das ist Feldlaufzeit, keine Weglaufzeit. Der Unterschied
ist physikalisch real:

\begin{itemize}[leftmargin=1.4em,itemsep=0.2em]
  \item Bei Weglaufzeit: B muss \emph{warten} bis das Signal ankommt.
    Vorher weiß B \emph{nichts}.
  \item Bei Feldlaufzeit: B wechselwirkt die ganze Zeit mit dem
    globalen Feld. Vorher reagiert B auf den ungestörten Feldzustand.
    Nachher reagiert B auf den durch die A-Messung modifizierten
    Feldzustand. Der Übergang erfolgt wenn die Wellenfront ankommt.
\end{itemize}

\textbf{[B]} Keine Information wurde von A nach B übertragen.
Die modifizierte Feldkonfiguration an B ist keine Botschaft von A ---
sie ist die natürliche Entwicklung des globalen Feldes nach der
A-Wechselwirkung. Da B den unmodifizierten Feldzustand sowieso
nicht kannte (Typ-II-Projektion, Phase $\theta$ unbekannt), kann B
auch aus dem modifizierten Feldzustand keine Information über die
A-Messung ziehen --- außer durch Vergleich mit A über einen
klassischen Kanal.

\subsection{Die $\xi$-Verzögerung als Feldphasenversatz}

\textbf{[K]} (A165) Die FFGFT-Vorhersage:
\[
  \Delta t_\xi = \xi \cdot \frac{d}{c} \approx 445\,\text{ns}
  \quad (d = 1000\,\text{km})
\]

\textbf{[S]} Diese Verzögerung ist \emph{keine} Weglaufzeit und
\emph{keine} zusätzliche Signalübertragung. Sie ist ein
\emph{Phasenversatz in der Feldwellenfront}, der durch die fraktale
Torusgeometrie entsteht.

Präzise: Das globale Feld auf dem $T^4$ ist nicht vollständig flach ---
auf irdischen Skalen gibt es einen Rest der Kompaktheit (A090). Dieser
Rest erzeugt eine geometrische Phasenverschiebung der Feldmoden.
Wenn die Wellenfront der Feldänderung sich von A aus ausbreitet,
trägt sie diesen Phasenversatz mit --- nicht weil sie einen Umweg
nimmt, sondern weil die Feldgeometrie auf dem Torus leicht von
der Euklidischen abweicht.

\textbf{[S]} Das bedeutet: Was Detektor B misst, ist die Fellamplitude
mit einem Phasenversatz $\xi \cdot d/c$ relativ zu dem, was B ohne
diesen geometrischen Effekt messen würde. Die 445\,ns sind die Zeit,
in der dieser Phasenversatz die Korrelationsstatistik verschiebt ---
nicht die Reisezeit eines Signals.

\textbf{[B]} Ob dieser Effekt messbar ist, hängt davon ab, ob man
die Korrelationsstatistik mit einer Zeitauflösung $\ll 445$\,ns
messen kann, relativ zu einer gemeinsamen Zeitreferenz Q.
Die Messung der Weglaufzeit (Signal A$\to$B) wäre konzeptuell falsch.
Die Messung der Feldlaufzeit (Änderung der Korrelationsstatistik
relativ zu $t_Q + d_B/c$) ist die richtige Frage.

\subsection{Warum es noch nie gemessen wurde}

Die bisherigen Bell-Experimente lassen sich in vier Generationen einteilen.
Keine davon hatte die Frage gestellt, die für die FFGFT-Vorhersage
entscheidend wäre.

\textbf{Generation~1: Nachweis der Bell-Verletzung.}
Aspect, Dalibard und Roger (1982) demonstrierten erstmals eine Bell-Verletzung
mit zeitveränderlichen Analysatoren \cite{Aspect1982}. Das Ziel war der
Nachweis, dass $S > 2$ möglich ist. Timing-Informationen wurden nur genutzt,
um Koinzidenzen zu identifizieren (Zeitfenster $\sim \mu$s), nicht um
die Korrelationsstruktur zeitlich aufzulösen.

\textbf{Generation~2: Schließung der Lokalitätslücke.}
Weihs et al.\ (1998) \cite{Weihs1998} schlossen die Lokalitätslücke durch
schnelle zufällige Wahl der Messrichtung ($\tau_\text{Entscheidung} \ll d/c$).
Die Zeitauflösung ($\sim 1$\,ns) war ausreichend für diesen Zweck, aber das
Experiment maß A gegen B --- nicht relativ zum Quellpunkt Q.

\textbf{Generation~3: Loophole-free Bell-Tests.}
Drei unabhängige Gruppen (Hensen et al.\ 2015 \cite{Hensen2015};
Giustina et al.\ 2015 \cite{Giustina2015}; Shalm et al.\ 2015 \cite{Shalm2015})
schlossen gleichzeitig Lokalitäts- und Detektionslücke.
Zeitauflösung: $\sim 1$\,ns, Abstände $d \sim 1$\,km.
Für $d = 1$\,km würde FFGFT $\Delta t \approx 0{,}44$\,ns vorhersagen ---
gerade noch unter der Zeitauflösung. Aber auch hier war die Fragestellung:
Verletzt $S$ die klassische Schranke? Nicht: Gibt es eine Verzögerung
in der Korrelationsstruktur?

\textbf{Generation~4: Kosmische Abstände.}
Yin et al.\ (2017) \cite{Yin2017} demonstrierten mit dem Micius-Satelliten
Bell-Verletzung über $\sim 1200$\,km. FFGFT sagt $\Delta t \approx 530$\,ns
voraus --- weit über der Nachweisgrenze von Atomuhren ($\sim 10$\,ns).
Aber Micius hatte keine präzise Quellpunkt-Referenz Q und keine drei
synchronisierten Uhren (Q, A, B) mit gemeinsamer Zeitbasis.

\textbf{Das strukturelle Problem.}
Alle vier Generationen messen dasselbe: ob die Korrelationsstatistik über
viele Ereignisse die klassische Schranke verletzt --- eine Ensemblefrage.
Die FFGFT-Vorhersage der 445-ns-Verzögerung ist dagegen eine Frage nach der
\emph{zeitlichen Struktur} der Feldmodifikation --- statistisch extrahiert,
aber konzeptuell verschieden.

Ein FFGFT-Test brauchte: (1) Präzise zeitgestempelten Quellpunkt Q.
(2) Drei synchronisierte Atomuhren Q, A, B.
(3) Statistische Analyse der Verzögerung an B relativ zu $t_Q + d_B/c$,
nicht relativ zu $t_A$.
(4) Ausreichende Ereigniszahl für $\Delta t \approx 445$\,ns.

Niemand hat dieses Experiment durchgeführt, weil die Fragestellung im
Standard-QM-Rahmen keinen Sinn ergibt: Dort propagiert nichts, der Kollaps
ist instantan. Die Frage ist nur unter einer Theorie sinnvoll, die eine
kausale Feldstruktur hinter den Korrelationen annimmt.

\begin{center}
\small
\begin{tabular}{lcccp{2.8cm}}
\toprule
Experiment & $d$ & $\Delta t_\text{FFGFT}$ & Auflösung & Fragestellung \\
\midrule
Aspect 1982 & $10$\,m & $\sim 0$ & $\mu$s & Bell-Verletzung \\
Weihs 1998 & $400$\,m & $0{,}18$\,ns & $\sim 1$\,ns & Lokalitätslücke \\
Hensen 2015 & $1{,}3$\,km & $0{,}58$\,ns & $\sim 1$\,ns & Loophole-free \\
Micius 2017 & $1200$\,km & $530$\,ns & $10$\,ns & Kosm.\ Distanz \\
\textbf{FFGFT} & $1000$\,km & $\mathbf{445}$\,ns & $\le 10$\,ns & $\xi$-Verzögerung? \\
\bottomrule
\end{tabular}
\end{center}

\begin{thebibliography}{9}
\small
\bibitem{Aspect1982}
Aspect, A., Dalibard, J., Roger, G. (1982).
\emph{Experimental test of Bell's inequalities using time-varying analyzers}.
Phys.\ Rev.\ Lett., 49(25), 1804--1807.
\bibitem{Weihs1998}
Weihs, G., Jennewein, T., Simon, C., Weinfurter, H., Zeilinger, A. (1998).
\emph{Violation of Bell's inequality under strict Einstein locality conditions}.
Phys.\ Rev.\ Lett., 81(23), 5039--5043.
\bibitem{Hensen2015}
Hensen, B. et al.\ (2015).
\emph{Loophole-free Bell inequality violation using electron spins
separated by 1.3 kilometres}. Nature, 526, 682--686.
\bibitem{Giustina2015}
Giustina, M. et al.\ (2015).
\emph{Significant-loophole-free test of Bell's theorem with entangled photons}.
Phys.\ Rev.\ Lett., 115(25), 250401.
\bibitem{Shalm2015}
Shalm, L. K. et al.\ (2015).
\emph{Strong loophole-free test of local realism}.
Phys.\ Rev.\ Lett., 115(25), 250402.
\bibitem{Yin2017}
Yin, J. et al.\ (2017).
\emph{Satellite-based entanglement distribution over 1200 kilometers}.
Science, 356(6343), 1140--1144.
\end{thebibliography}

\section{Verwendung von KI-Werkzeugen}

Dieses Dokument ist unter Verwendung KI-gestützter Werkzeuge entstanden. Die
Fragestellungen, die Setzungen und alle inhaltlichen Entscheidungen stammen vom
Autor; die Werkzeuge formulieren aus, rechnen nach und erstellen Prüfskripte.
Die Verantwortung für Inhalt und Fehler liegt vollständig beim Autor. Der
ausführliche Wortlaut steht in A010.


